\subsection{مبانی نظری و فرمول‌بندی ریاضی}
فرض کنید $\mathbf{W} \in \mathbb{R}^{d_{\text{out}} \times d_{\text{in}}}$ ماتریس وزن یک لایه خطی متراکم باشد که می‌توان آن را به صورت مجموعه‌ای از بردارهای ردیفی (نورون‌های خروجی) نمایش داد:
\begin{equation}
\mathbf{W} = \begin{bmatrix} \mathbf{w}_1^\top \\ \mathbf{w}_2^\top \\ \vdots \\ \mathbf{w}_{d_{\text{out}}}^\top \end{bmatrix}, \quad \mathbf{w}_i \in \mathbb{R}^{d_{\text{in}}}
\end{equation}

الگوریتم \lr{ToMoE} با اجرای خوشه‌بندی \lr{K-Means} با دقت محاسباتی \lr{FP32} (یا \lr{MiniBatchKMeans} برای لایه‌های بسیار بزرگ با $d_{\text{out}} > 2048$) نورون‌ها را به $E$ خوشه متمایز $\mathcal{C}_1, \mathcal{C}_2, \dots, \mathcal{C}_E$ با ابعاد ترجیحاً برابر افراز می‌نماید:
\begin{equation}
\arg\min_{\{\mathcal{C}_e\}_{e=1}^E} \sum_{e=1}^E \sum_{\mathbf{w}_i \in \mathcal{C}_e} \|\mathbf{w}_i - \boldsymbol{\mu}_e\|_2^2
\end{equation}
که در آن $\boldsymbol{\mu}_e$ مرکز خوشه $e$-ام است. در صورت تولید خوشه‌های خالی یا نابرابری شدید، مکانیزم پشتیبان افراز مساوی و معین (\lr{Deterministic Equal Division}) فعال می‌شود.

هر خوشه متناظر با یک خبره $\mathbf{W}^{(e)} \in \mathbb{R}^{d_e \times d_{\text{in}}}$ است. برای هر بردار ورودی $\mathbf{x} \in \mathbb{R}^{d_{\text{in}}}$، شبکه مسیریاب با استفاده از نگاشت خطی $\mathbf{W}_{\text{gate}} \in \mathbb{R}^{E \times d_{\text{in}}}$ امتیازهای هدایت را در فضای \lr{FP32} محاسبه می‌کند:
\begin{equation}
\mathbf{h} = \mathbf{x} \mathbf{W}_{\text{gate}}^\top + \mathbf{b}_{\text{gate}}
\end{equation}
سپس، عملگر انتخاب برترین شاخص‌ها (\lr{Top-k}) وزن‌های ترکیب را تعیین می‌نماید:
\begin{equation}
\mathcal{T} = \operatorname{Top-k}(\mathbf{h}, k), \quad g_e = \begin{cases} \frac{\exp(h_e)}{\sum_{j \in \mathcal{T}} \exp(h_j)} & \text{if } e \in \mathcal{T} \\ 0 & \text{otherwise} \end{cases}
\end{equation}

خروجی موثر نهایی از پراکنده‌سازی و بازنشانی خروجی خبره‌های فعال بر روی بردار کامل بعد خروجی بازسازی می‌شود:
\begin{equation}
\mathbf{y} = \sum_{e \in \mathcal{T}} g_e \cdot \operatorname{Scatter}\left(\mathbf{x} (\mathbf{W}^{(e)})^\top + \mathbf{b}^{(e)}, \mathcal{I}_e\right)
\end{equation}
که در آن $\mathcal{I}_e$ اندیس نورون‌های متعلق به خوشه $e$ در ماتریس خروجی اصلی است.


\paragraph{تحلیل طیفی و نامساوی‌های پایداری خطا}
برای بررسی میزان دقت تقریب ماتریس وزن یا نگاشت توجه در ماژول \lr{tomoe\_compressor}، تفاضل ماتریس فاکتورگیری‌شده از نگاشت متراکم اصلی بر اساس نامساوی هافدینگ-فون‌نیومان ارزیابی می‌گردد. فرض کنید $\mathbf{W} \in \mathbb{R}^{d_{\text{in}} \times d_{\text{out}}}$ ماتریس وزن هدف باشد که توسط تقریب کم‌مرتبه یا پراکنده با رتبه موثر $r$ جایگزین شده است. کران بالای خطای بازسازی بر مبنای نرم فروبنیوس به صورت زیر تصریح می‌شود:
\begin{equation}
\label{eq:spectral_bound_tomoe_compressor}
\| \mathbf{W} - \mathbf{W}_{r} \|_F = \sqrt{\sum_{i=r+1}^{\min(d_{\text{in}}, d_{\text{out}})} \sigma_i^2(\mathbf{W})} \le \epsilon \| \mathbf{W} \|_F
\end{equation}
که در آن $\sigma_i(\mathbf{W})$ بیانگر $i$-امین مقدار منفرد مرتب‌شده ماتریس به صورت نزولی است. نسبت حفظ انرژی طیفی $\eta_r = \frac{\sum_{i=1}^r \sigma_i^2}{\sum_{i=1}^{d} \sigma_i^2}$ تضمین می‌کند که با انتخاب آستانه $\tau \ge 0.90$، بیش از نود درصد از پراکندگی واریانس داده‌ها حفظ شده و افت سیگنال در انتشار گرادیان رو به عقب به حداقل ممکن کاهش می‌یابد.

\paragraph{تحلیل مجانبی و پیچیدگی اشتراک حافظه سخت‌افزاری}
از منظر پیچیدگی محاسباتی و انتقال داده‌ها در سلسومراتب حافظه (حافظه سراسری \lr{HBM} به حافظه محلی \lr{SRAM})، عملیات ضرب ماتریسی در این مکانیزم نیازمند تخصیص بافر موقت با ابعاد فاکتورگیری‌شده است. برای پردازش دنباله‌ای به طول $N$ با بعد مخفی $D$، بار محاسباتی هسته از فرم درجه دو به رابطه زیر تقلیل می‌یابد:
\begin{equation}
\label{eq:complexity_tomoe_compressor}
\mathcal{C}_{\text{tomoe\_compressor}} = \mathcal{O}\left( B \cdot N \cdot D \cdot r + B \cdot H \cdot N \cdot d_k \right) \ll \mathcal{O}\left( B \cdot H \cdot N^2 \cdot d_k \right)
\end{equation}
این بهبود مرتبه پیچیدگی، با کاهش چشمگیر حجم تبادل تانسورها در گذرگاه‌های ارتباطی کارت گرافیک، مانع از اشباع پهنای باند حافظه (پدیده \lr{Memory-Bound}) در فرآیند استنتاج طولانی‌مدت می‌گردد.
